\chapter{进程、线程与调度}

\section{原理}

操作系统的第一个核心职责是分配 CPU 时间。硬件只有有限核心，用户却希望多个程序同时推进。进程提供资源边界，线程提供执行边界，调度器在可运行实体之间选择下一个执行者。这里的关键不是“轮流执行”四个字，而是每个执行实体必须处在可解释的状态里。

先从一个很小的服务器看问题：

\begin{lstlisting}
while true:
  c = accept(listen_fd)
  data = read(c)
  write(log_fd, data)
  write(c, "ok")
\end{lstlisting}

如果这段代码只有一个执行流，问题马上出现：没有客户端连接时，\code{accept} 会等待；客户端发数据慢时，\code{read} 会等待；磁盘慢时，写日志会等待。CPU 在这些等待期间其实可以运行别的请求，但硬件不会自动知道“这个程序在等网络，那个程序可以继续算”。内核必须把“能继续执行”和“正在等待某个条件”区分成状态，再让调度器只选择能继续执行的实体。

最朴素的错误调度器会这样写：

\begin{lstlisting}
for task in all_tasks:
  run(task)
\end{lstlisting}

它的问题不是代码短，而是没有状态。一个正在等磁盘的任务被运行也没有意义；一个已经退出的任务不该再运行；一个刚被网卡中断唤醒的任务应该重新进入候选集合；一个持有锁的内核路径不能在任意位置被切走。于是调度章真正要回答的不是“怎么排序”，而是四件事：任务状态怎么表示，等待条件怎么保存，唤醒怎样回到 runqueue，上下文怎样被保存和恢复。

进程不是一个正在运行的函数。它至少包含地址空间、文件描述符表、信号处理状态、凭据、资源限制、子进程关系和一个或多个线程。线程是内核可调度实体，拥有寄存器上下文、栈、调度状态和等待原因；同一进程内的线程共享地址空间和许多进程资源，所以它们能高效共享数据，也更容易互相破坏。

这里第一次出现两个容易混的词。进程是资源容器：它回答“这个程序能看见哪些地址、打开了哪些 fd、以什么身份运行”。线程是执行实体：它回答“当前 RIP/RSP/寄存器在哪里，能不能被调度到 CPU 上”。一个单线程进程看起来像两者合一；多线程进程把区别暴露出来：多个线程共享同一个地址空间和 fd 表，但每个线程有自己的栈、寄存器现场和等待状态。

\begin{longtable}{p{0.2\linewidth}p{0.34\linewidth}p{0.34\linewidth}}
\toprule
概念 & 回答的问题 & 常见误解 \\
\midrule
进程 & 资源和权限边界是什么 & 误以为进程就是正在跑的那条函数调用链 \\
线程 & 哪个执行现场可以放到 CPU 上 & 误以为线程只是一段代码，而不是保存现场的对象 \\
runqueue & 哪些线程现在可以运行 & 误把所有未退出线程都放进候选集合 \\
wait queue & 哪些线程在等同一个条件 & 误以为睡眠线程会自动被 CPU 记住 \\
上下文切换 & 保存当前现场，恢复另一个现场 & 误以为只改一个 current 指针就够 \\
\bottomrule
\end{longtable}

进程生命周期可以简化为 created、runnable、running、blocked、zombie、reaped。线程状态可以简化为 runnable、running、sleeping、stopped、terminated。真实内核状态更多，但第一版必须先守住这些边界：blocked 线程不能被调度执行，terminated 线程不能重新运行，zombie 进程已经退出但仍等待父进程回收退出状态。

调度策略必须回答：谁处于 runnable，谁在 blocked，时间片用完后怎么办，优先级是否影响选择，任务退出后资源如何回收，唤醒是否会造成饥饿。一个可解释的调度器必须把这些状态外化，而不是只返回一个线程编号。

上下文切换是调度的成本边界。切换时内核要保存当前线程的寄存器、栈指针和调度账本，选择下一个线程，恢复它的上下文，并可能改变地址空间和 TLB 状态。线程切换不一定改变地址空间，进程切换通常会带来更重的内存管理成本。高并发程序如果制造过多 runnable 线程，就会把时间花在 runqueue 等待、抢占和上下文切换上。

等待是调度的另一半。一个线程不运行，可能是被抢占，也可能是主动睡眠，可能等锁、等 I/O、等 page fault、等子进程、等信号、等定时器。分析性能和死锁时，必须区分 on-CPU 时间和 off-CPU 时间。只看 CPU 使用率会误导：CPU 低可能是线程都阻塞了，CPU 高可能是大量线程在抢锁或忙等。

后续章节的很多机制都会回到这套状态机。缺页时，线程可能睡在“等待磁盘把文件页读入”的队列上；socket 没数据时，线程睡在 socket receive wait queue 上；\code{waitpid} 睡在子进程退出事件上；\code{fsync} 睡在写回 completion 上。所谓操作系统“同时处理很多事”，本质是把不能前进的执行流从 CPU 上拿下来，并保证条件满足时能找回它。

\section{实践}

纸上实践先画三张表。第一张是进程表，列出 PID、父 PID、状态、退出码、地址空间和 fd 表引用。第二张是线程表，列出 TID、所属 PID、状态、等待原因、优先级、最近运行 CPU 和已消耗 tick。第三张是 runqueue，列出当前可运行线程的顺序和调度账本。

\begin{longtable}{@{}p{0.08\linewidth}p{0.17\linewidth}p{0.2\linewidth}p{0.14\linewidth}p{0.25\linewidth}@{}}
\toprule
tick & 事件 & runqueue & 运行线程 & 状态变化 \\
\midrule
0 & 创建 A/B & A,B & A & A running \\
1 & 时间片结束 & B,A & B & A runnable, B running \\
2 & B 发起 I/O & A & A & B blocked(io) \\
3 & I/O 完成 & A,B & A & B runnable \\
4 & A exit & B & B & A zombie \\
\bottomrule
\end{longtable}

这张表比一句“轮转调度”更有价值，因为它同时呈现了状态、队列和事件。后续任何复杂策略都可以在这张表上增加字段，例如虚拟运行时间、deadline、优先级、CPU affinity、NUMA node 或 cgroup quota。

实践报告必须区分三个时间：wall time、runnable wait 和 on-CPU time。wall time 是请求从开始到结束的总时间；runnable wait 是线程已经可运行但还没拿到 CPU 的时间；on-CPU time 是真正执行指令的时间。把三者混在一起，会把调度问题误判成算法问题。

\section{算法与实现分析}

\subsection{FCFS：最简单也最容易产生长等待}

先来最简单的调度算法：first-come, first-served。任务按进入 runnable 队列的顺序执行，一个任务运行到阻塞、退出或主动让出 CPU 后，才轮到下一个任务。

\begin{lstlisting}
while true:
  task = runqueue.pop_front()
  run_until_block_or_exit(task)
  if task.state == RUNNABLE:
    runqueue.push_back(task)
\end{lstlisting}

FCFS 的实现成本很低：入队 \code{O(1)}，出队 \code{O(1)}。它的问题是 convoy effect。一个长 CPU 任务排在前面，后面的短交互任务会长时间等待。若系统只跑批处理，FCFS 可以作为 reference；若系统需要交互响应，它通常不合格。

\subsection{Round-robin：用时间片限制独占}

Round-robin 在 FCFS 上增加时间片。每个 runnable 任务最多运行一个 quantum，时间片到期后被放回队尾。它解决了“一个任务长期霸占 CPU”的问题，但引入了时间片选择：时间片太大，交互延迟高；时间片太小，上下文切换成本高。

\begin{lstlisting}
on timer_tick:
  current.used_ticks += 1
  if current.used_ticks == current.quantum:
    current.used_ticks = 0
    current.state = RUNNABLE
    runqueue.push_back(current)
    current = runqueue.pop_front()
\end{lstlisting}

假设上下文切换成本是 S，时间片是 Q，CPU 时间被切换消耗的比例大约是 \code{S / (S + Q)}。Q 越小，响应越好，但浪费越多。这个公式不是精确模型，却能解释为什么时间片不能无限小。

\subsection{优先级调度和饥饿}

优先级调度把 runqueue 拆成多个队列，总是选择最高优先级的 runnable 任务。它适合表达“控制任务比日志任务重要”，但如果高优先级任务持续可运行，低优先级任务可能永远得不到 CPU。

\begin{lstlisting}
for priority from HIGH to LOW:
  if !queue[priority].empty():
    return queue[priority].pop_front()
return idle_task
\end{lstlisting}

解决饥饿的常见方法是 aging：等待越久，优先级逐步提高。aging 的实现要小心，不能每个 tick 扫描所有任务，否则任务多时成本过高。可以在任务入队时记录等待起点，在选择或周期性维护时计算有效优先级。

\subsection{MLFQ：用反馈区分交互和批处理}

Multilevel feedback queue 使用多个优先级队列，新任务先进入高优先级；如果任务用完整个时间片，说明它偏 CPU 密集，下次降级；如果任务很快阻塞或让出 CPU，说明它可能是交互任务，保留较高优先级。MLFQ 不需要提前知道任务类型，而是从行为反馈中调整。

\begin{lstlisting}
on_task_created(task):
  task.level = 0
  queues[0].push_back(task)

on_quantum_expired(task):
  task.level = min(task.level + 1, MAX_LEVEL)
  queues[task.level].push_back(task)

on_task_blocked_before_quantum(task):
  queues[task.level].push_back_when_woken(task)
\end{lstlisting}

MLFQ 的关键参数包括层数、每层时间片、降级规则、周期性 boost。boost 用来防止低层任务永久饥饿：每隔一段时间，把所有 runnable 任务提升回高层。实现上要记录任务当前层级、已用时间、最近 boost epoch。

\subsection{公平调度：从 CFS 模型到 EEVDF}

公平调度不只问“谁先来”，而是问每个 runnable 任务已经得到多少 CPU 服务。经典 CFS 模型给每个任务维护虚拟运行时间 \code{vruntime}：任务实际运行越久，\code{vruntime} 越大；权重越高，虚拟时间增长越慢。下面的伪代码是用于理解 CFS 核心思想的简化模型，不是当前 Linux 调度器的完整实现。

\begin{lstlisting}
on_task_run(task, delta_exec):
  task.vruntime += delta_exec * NICE_0_WEIGHT / task.weight

pick_fair(rbtree):
  return leftmost_node_by_vruntime(rbtree)
\end{lstlisting}

若用红黑树保存 runnable 任务，插入、删除大约是 \code{O(log n)}，取缓存的最左节点可接近 \code{O(1)}。这个模型解释了为什么 nice 值不是简单固定优先级：它影响获得 CPU 的比例，而不是让低优先级任务永远不运行。

Linux 从 6.6 起逐步把公平类的选择策略迁移到 EEVDF。EEVDF 仍使用虚拟时间核算服务量，但先用 lag 判断任务是否有资格运行，再在有资格的任务中选择虚拟截止期最早者。阅读 \filepath{kernel/sched/fair.c} 时，应把“CFS 的最小 \code{vruntime} 教学模型”和“当前 EEVDF 的 eligibility + virtual deadline 决策”分开。

\subsection{deadline 调度与准入控制}

实时或准实时任务还会声明运行时间和截止时间。例如任务说“每 20ms 需要 3ms CPU”。调度器若盲目接受所有请求，总利用率超过 100 后必然 miss deadline。因此 deadline 调度需要准入控制。

\begin{lstlisting}
admit(task):
  new_util = total_runtime_over_period + task.runtime / task.period
  if new_util > CPU_CAPACITY:
    return -EBUSY
  add_to_deadline_class(task)
  return OK

pick_deadline():
  return task with earliest absolute_deadline
\end{lstlisting}

EDF 在单 CPU 理想模型中有很强性质，但真实系统还有抢占成本、临界区、缓存失效、中断关闭和共享资源阻塞。理论可调度不等于工程一定按时，测试必须记录 deadline miss、最大关抢占时间和锁等待。

\subsection{调度指标}

调度器优化必须有指标。常用指标包括平均响应时间、p95/p99 延迟、吞吐、上下文切换次数、公平误差、迁移次数、cache miss、deadline miss 和能耗。只报告“更快”没有意义。

\begin{longtable}{p{0.22\linewidth}p{0.34\linewidth}p{0.32\linewidth}}
\toprule
指标 & 说明 & 常见误判 \\
\midrule
on-CPU time & 真正执行指令的时间 & 把等待 I/O 误认为算法慢 \\
runnable wait & 可运行但没拿到 CPU 的时间 & 把调度拥塞误认为锁死 \\
context switches & 切换次数和成本 & 只看次数，不看每次成本 \\
migration count & 跨 CPU 迁移 & 均衡改善但缓存局部性变差 \\
fairness error & 实得 CPU 与权重目标差距 & 平均公平掩盖尾延迟 \\
\bottomrule
\end{longtable}

\subsection{真实系统为什么更复杂}

真实通用 OS 还要处理多核、NUMA、CPU affinity、实时任务、cgroup 配额、负载均衡和能耗。多核调度不是简单地把一个队列给所有 CPU，因为全局队列会产生锁竞争；每 CPU runqueue 能降低竞争，但需要负载均衡。迁移任务可以平衡负载，却会丢失 cache/TLB 热状态。调度器必须在公平、吞吐、延迟、局部性和能耗之间折中。

\subsection{task 不是一个正在运行的函数}

硬件正在执行的是寄存器、RIP、RSP、RFLAGS 和页表上下文。内核把这组执行状态包装成 task，再让 task 指向地址空间、文件表、信号状态、凭据、namespace 和 cgroup。task 是可调度实体，不是进程所有资源的简单大盒子。

\begin{longtable}{p{0.20\linewidth}p{0.34\linewidth}p{0.34\linewidth}}
\toprule
字段类别 & 保存什么 & 为什么重要 \\
\midrule
执行现场 & 内核栈、保存寄存器、返回路径 & 被切走后还能从同一点继续 \\
调度状态 & state、policy、prio、vruntime、on\_rq & 判断能否运行、怎样排序、是否已入队 \\
地址空间 & \code{mm}、\code{active\_mm} & 进程切换可能要切页表和 TLB \\
资源指针 & files、fs、signal、sighand & fork、exec、exit 要复制、共享或释放 \\
安全身份 & cred、namespace、cgroup & 权限、限额和可见性都依赖它 \\
亲和性 & allowed CPUs、last CPU、NUMA 统计 & 负载均衡不能只看队列长度 \\
\bottomrule
\end{longtable}

\code{state} 和 \code{on\_rq} 要分开理解。\code{state} 说明任务是 runnable、sleeping、stopped 还是 exiting；\code{on\_rq} 说明它是否已经挂到某个 runqueue。一个 runnable 任务没有入队，会永远等不到 CPU；一个 sleeping 任务留在 runqueue，会被错误调度。调度器的大量锁和状态检查，就是为了让这两本账一致。

\subsection{阻塞系统调用怎样离开 CPU}

阻塞不是“函数停住”。内核要把当前任务从可运行集合移走，把它挂到等待队列，并保证事件到达时能找回它。以等待 pipe 或 socket 可读为例：

\begin{lstlisting}
blocking_read(obj):
  lock(obj.lock)
  while !data_available(obj):
    prepare_to_wait(obj.read_waitq, current)
    current.state = TASK_INTERRUPTIBLE
    unlock(obj.lock)
    schedule()
    lock(obj.lock)
    finish_wait(obj.read_waitq, current)
    if signal_pending(current):
      unlock(obj.lock)
      return -EINTR
  copy_data_to_user()
  unlock(obj.lock)
  return bytes
\end{lstlisting}

检查谓词和入等待队列必须受同一把锁保护。若先检查发现没有数据，还没入队时生产者写入并唤醒，唤醒会丢失；随后当前任务睡下，就形成 lost wakeup。谓词循环也不是模板写法：任务可能被信号唤醒，可能被虚假唤醒，也可能多个等待者竞争同一份数据。

把这个错误具体走一遍，问题会更清楚。假设线程 A 正在 \code{read(pipefd)}，线程 B 稍后 \code{write(pipefd)}。错误实现若按“先看 pipe 为空，再准备睡眠”两步分开做，中间没有锁保护，就会出现下面的时序：

\begin{longtable}{p{0.15\linewidth}p{0.38\linewidth}p{0.35\linewidth}}
\toprule
时刻 & 线程 A & 线程 B 或内核状态 \\
\midrule
1 & A 检查 pipe，发现没有数据 & wait queue 仍为空 \\
2 & A 还没把自己挂入 wait queue & B 写入数据，调用 wakeup，但没有等待者可唤醒 \\
3 & A 把自己设为 sleeping 并调用 \code{schedule} & pipe 已经有数据，但没有新的 wakeup 会来 \\
4 & CPU 去运行别的任务 & A 可能永久睡眠，除非信号、超时或其他写入再次触发事件 \\
\bottomrule
\end{longtable}

正确实现把“检查条件、登记等待者、改变 task 状态”放在同一个临界区里。这样生产者要么在 A 入队前看见锁被持有，等 A 完成睡眠准备后再写入并唤醒；要么在 A 检查条件前已经写入数据，A 重新检查时发现有数据而不会睡下。这里的锁不是为了保护几行代码好看，而是在保护一个跨 CPU 的事实：若条件已经满足，等待者不能睡丢；若等待者已经睡下，满足条件的一方必须能找到它。

还要注意，\code{schedule()} 返回不表示数据一定已经到手。A 可能因为信号返回，可能因为超时返回，可能被广播唤醒后发现另一个线程先取走了数据。因此睡眠后必须重新拿锁、重新检查谓词。许多初学者会把 \code{wake\_up} 当成“把函数从睡眠点继续执行”，这不准确。真实过程是：wakeup 修改任务状态并把它放回 runqueue；真正执行要等调度器选择它；继续执行后还要重新证明条件仍然成立。

\subsection{唤醒怎样进入 runqueue}

唤醒不等于立即运行。唤醒方通常只把等待任务改回 runnable，放到目标 CPU 的 runqueue；是否抢占当前任务，要看调度类、优先级、抢占状态和 CPU 位置。

\begin{lstlisting}
wake_up(waitq):
  lock(waitq.lock)
  for task in waitq:
    if task.wait_condition_is_satisfied:
      remove_from_waitq(task)
      task.state = TASK_RUNNING
      enqueue_task(select_runqueue(task), task)
      if task_should_preempt_current(task):
        set_need_resched(target_cpu)
  unlock(waitq.lock)
\end{lstlisting}

若目标 CPU 正在用户态运行，\code{need\_resched} 常在返回用户态、定时器中断返回或显式调度点处理。若目标 CPU 是另一个核心，内核可能发 IPI 让它尽快重新调度。这样看，wakeup 的成本包括等待队列锁、runqueue 锁、跨 CPU 通知和 cache line 迁移。

再用一个具体的 socket 接收案例串起来。线程 A 在 CPU0 上调用 \code{recv}，接收队列为空，于是它把等待原因记成“等 socket receive queue 非空”，挂到 socket 的等待队列，然后从 CPU0 的 runqueue 中消失。此时 A 没有“暂停在 CPU 里”，它只是一个 task 对象，保存了内核栈、寄存器现场、等待队列链接和状态。CPU0 可以去运行线程 C。

随后网卡在 CPU1 上产生中断。驱动或 NAPI poll 把包交给协议栈，协议栈找到目标 socket，把数据放进 receive queue，然后调用 wakeup。wakeup 要做的不是直接跳回线程 A 的用户栈，而是把 A 从 socket wait queue 摘掉，改成 runnable，选择一个目标 runqueue。若 A 上次在 CPU0 运行且 cache/TLB 仍有局部性，调度器可能倾向放回 CPU0；若 CPU0 很忙而 CPU1 空闲，也可能迁移。若迁移，A 将来恢复时会付出 cache 冷启动成本。

\begin{longtable}{p{0.18\linewidth}p{0.36\linewidth}p{0.34\linewidth}}
\toprule
阶段 & 状态变化 & 读者要抓住的边界 \\
\midrule
\code{recv} 阻塞 & task 从 runqueue 移到 socket wait queue & 它不再竞争 CPU，但内核必须记住等待条件 \\
包到达 & 数据进入 socket receive queue & 条件变为真，但等待线程还没有运行 \\
wakeup & task 变成 runnable，进入某个 runqueue & runnable 只表示有资格运行，不表示已经运行 \\
调度选择 & CPU 在安全点切换到该 task & 可能立即抢占，也可能等当前任务返回或时间片结束 \\
重新检查 & \code{recv} 继续后再次检查 receive queue & 多等待者、信号和竞态要求谓词循环 \\
\bottomrule
\end{longtable}

这个案例解释了为什么排查“请求卡住”不能只看线程是不是 sleeping。若线程还在 wait queue，问题可能是事件没发生、唤醒路径漏了、等待条件写错；若线程已经 runnable 但很久没运行，问题在调度拥塞；若线程运行后又睡回去，可能是虚假唤醒、条件被别的线程抢走，或协议状态不满足。状态分清楚，才有办法定位。

\subsection{抢占点和不可抢占区}

定时器中断可以打断用户态程序，但内核不能在任意位置无条件切走当前任务。持有自旋锁、关闭中断、正在修改 runqueue、正在操作每 CPU 数据时，随意抢占会破坏不变量。内核用 \code{preempt\_count}、中断状态和锁规则标记这些区域。

\begin{lstlisting}
timer_interrupt:
  account_runtime(current)
  if current.timeslice_expired:
    set_need_resched(current)
  if returning_to_user and need_resched:
    schedule()
  else if kernel_preemptible and preempt_count == 0:
    schedule()
\end{lstlisting}

实时系统关心的不是平均时间片，而是最长不可抢占区、最长关中断区和最高优先级任务被低优先级持锁路径阻塞的时间。一个系统平均响应很好，也可能因为某条持锁路径过长而出现不可接受的尾延迟。

\subsection{fork、exec、wait 和调度器的连接}

\code{fork} 创建的新任务不能在半初始化状态进入 runqueue。内核必须先准备内核栈、保存的用户寄存器、调度实体、地址空间引用、文件表引用、信号状态和 PID，再让调度器看到它。子进程第一次运行时不是从程序开头执行，而是从系统调用返回路径回到用户态，只是返回值寄存器被设置为 0。

\begin{lstlisting}
fork scheduling boundary:
  allocate child task
  prepare child saved registers
  copy or share mm/files/signal/cred
  allocate pid and link parent-child relation
  enqueue child as runnable
  parent returns child pid
  child later returns 0 from copied frame
\end{lstlisting}

\code{exec} 成功后仍是同一个 task 被调度，但地址空间、用户栈、入口 RIP、部分信号处理和 close-on-exec fd 改变。\code{wait} 要求已退出子进程保留 zombie 壳，直到父进程读取退出码。调度器只负责不再运行 zombie；进程管理还要保证退出状态只被回收一次，最后一个引用释放 task。

\subsection{源码阅读入口}

进程创建看 \filepath{kernel/fork.c} 的 \code{copy\_process}；退出和回收看 \filepath{kernel/exit.c} 的 \code{do\_exit}、\code{release\_task}；调度主干看 \filepath{kernel/sched/core.c} 的 \code{\_\_schedule}、\code{try\_to\_wake\_up}；公平类看 \filepath{kernel/sched/fair.c} 的 enqueue/dequeue、vruntime 和负载均衡。阅读时把每一步标成四类动作：改 task 状态、改 runqueue、改资源引用、通知等待者。

\subsection{本章压缩进来的并发与锁}

调度章不能只到“谁获得 CPU”为止。只要两个线程可能同时碰同一份状态，就必须说明同步对象、等待对象和失败恢复。锁、futex、信号、IPC 和死锁原本可以展开成很多章，但主线里先抓住一个统一问题：某个线程为什么不能继续，它在等谁，谁有资格改变条件，条件改变后谁负责唤醒。

\begin{longtable}{p{0.20\linewidth}p{0.36\linewidth}p{0.32\linewidth}}
\toprule
机制 & 它保存的等待事实 & 常见错误 \\
\midrule
mutex & 持有者是谁，等待者排在哪个队列 & 忘记在循环中复查条件，或异常路径漏解锁 \\
futex & 用户态字值和内核等待队列的绑定 & 只看内核队列，不验证用户态谓词是否仍成立 \\
condition variable & 等待哪个共享状态变真 & 把 wakeup 当成状态已经满足 \\
signal & 异步事件和可中断睡眠的交汇点 & 忽略 \code{EINTR}，把部分完成误当失败 \\
pipe/socket & 缓冲区从空到非空、从满到非满 & 读写端关闭后仍然等待不可能发生的事件 \\
\bottomrule
\end{longtable}

锁的规则也必须落到调度状态。自旋锁适合短临界区，持锁时不能睡眠；mutex 可睡眠，适合较长临界区；读写锁提高读多写少场景的并发，但写者饥饿和升级死锁要单独处理；RCU 让读侧几乎不阻塞，但释放对象必须等 grace period，不能在最后一个指针消失时立即 \code{free}。这些机制的差别，不是 API 名字，而是“持有期间能不能调度、等待者在哪里、对象何时真正可释放”。

\begin{lstlisting}
safe_wait(lock, waitq, condition):
  lock(lock)
  while !condition():
    prepare_to_wait(waitq, current)
    unlock(lock)
    schedule()
    lock(lock)
  finish_wait(waitq, current)
  unlock(lock)
\end{lstlisting}

这段伪代码压缩了等待路径的核心不变量：检查条件和入队必须受同一把锁保护；睡眠前要释放保护业务状态的锁；醒来后要重新加锁并复查条件；退出等待时要清理 wait queue 链接。若先检查条件再入队，中间事件可能已经发生并唤醒了空队列，线程随后睡下就永远没人叫醒。若醒来后不复查，多等待者会把同一个资源消费两次。

死锁分析也要用对象图，不要只说“锁顺序错了”。把每个线程正在持有的锁、正在等待的锁、是否可被信号打断、是否持有不可睡眠锁都列出来。如果图中出现环，而且环上没有超时、取消或外部释放路径，系统就无法自行前进。lockdep 一类工具本质上是在运行时收集锁类顺序，提前发现可能形成环的获取顺序。

\begin{rubricbox}
调度与并发章的最低验收：给任一阻塞点，都能写出等待谓词、等待队列、唤醒者、是否可中断、被唤醒后的复查动作、持锁规则和失败清理路径。缺一项，代码即使在正常路径能跑，也没有可证明的并发语义。
\end{rubricbox}

\section{测试}

调度测试至少覆盖三类情况：多个 runnable 进程应轮转推进；blocked 进程不应获得 tick；terminated 进程不应重新出现。还要检查没有 runnable 进程时，调度器返回 idle，而不是访问空队列。

\begin{itemize}
\tightlist
\item fork 后父子进程关系可解释，父进程能回收子进程退出状态。
\item exec 后地址空间替换，但 PID、部分 fd 和进程身份按规则保留。
\item 线程阻塞后从 runqueue 消失，唤醒后回到 runnable 集合。
\item 时间片结束只改变调度状态，不应破坏进程资源。
\item 没有 runnable 线程时进入 idle，而不是忙等消耗 CPU。
\item 优先级或 aging 不得让低优先级任务永久饥饿，除非策略明确允许。
\end{itemize}

后续可以加入优先级、deadline、affinity、负载均衡和 cgroup 配额，但必须保留基础回归。任何复杂策略都不能让 blocked 或 terminated 任务被错误调度。

\inputdetail{chapters/ch05-process-scheduler-implementation-deep}
